%% ma: L11-B14-0014
%% lop: 11
%% chuong: 4
%% bai: 14
%% dang: 11.14.1
%% loai: TN4
%% muc: NB
%% dap_an: "A"
%% nguon: "tai-lieu-hoc-tap-mon-toan-lop-11-theo-chuong-trinh-giao-duc-pho-thong-2018.pdf (Lớp Toán Thầy Duy, Phạm Lê Duy, THPT Chu Văn An), Chương 4 Bài 5, Phần 1, Câu 9 (tr.528)"
%% kien_thuc: "Lớp 11 Bài 14 (phép chiếu song song)"
%% trang_thai: cho-duyet
%% ly_do: "Dạng mới đề xuất (11.14.x), chờ giáo viên duyệt dạng; tài liệu gốc không có lời giải mẫu, đáp án do kho tự giải và kiểm bằng sympy."
\begin{ex}
Cho hình lăng trụ $ABC.A^{\prime}B^{\prime}C^{\prime}$. Gọi $G$ là trọng tâm của tam giác $ABC$. Qua phép chiếu song song đường thẳng $AA^{\prime}$ mặt phẳng chiếu là $(A^{\prime}B^{\prime}C^{\prime})$ biến $G$ thành $G^{\prime}$. Tìm mệnh đề đúng?
\choice{\True $G^{\prime}$ là trọng tâm tam giác $A^{\prime}B^{\prime}C^{\prime}$.}{$G^{\prime}$ là trung điểm của $A^{\prime}B^{\prime}$.}{$G^{\prime}$ là trực tâm của tam giác $A^{\prime}B^{\prime}C^{\prime}$.}{$G^{\prime}$ là trung điểm của $B^{\prime}C^{\prime}$.}
\loigiai{\sloppy Phương chiếu $AA^{\prime}\parallel BB^{\prime}\parallel CC^{\prime}$ nên $A,B,C\mapsto A^{\prime},B^{\prime},C^{\prime}$. Phép chiếu song song bảo toàn tỉ số độ dài trên cùng đường thẳng nên trọng tâm $G$ của $\triangle ABC$ biến thành trọng tâm $G^{\prime}$ của $\triangle A^{\prime}B^{\prime}C^{\prime}$.}
\end{ex}
